\section{Data and exploratory analysis}
\label{sec:data_eda}

We use two related but distinct evaluation settings. The first is a cross-well
clay-volume benchmark using noisy 30 dB logs, where the training wells are
F02-1, F03-2, and F06-1 and the held-out test well is F03-4. The second is a
single-well F03-4 porosity benchmark using only gamma ray as the dense input
channel. Table~\ref{tab:eda_dataset_summary} and
Fig.~\ref{fig:eda_target_summary} summarize the target distributions used by
both settings.

\begin{table}[htbp]
  \centering
  \caption{Dataset summary used to motivate the two experimental settings.}
  \label{tab:eda_dataset_summary}
  \small
  \resizebox{\linewidth}{!}{\begin{tabular}{lllrrrrrrr}
\toprule
dataset & well & target & n\_rows & depth\_min & depth\_max & target\_mean & target\_std & target\_min & target\_max \\
\midrule
cross\_train & F02-1 & vc & 165 & 1149.7056 & 1174.6992 & 0.1215 & 0.0407 & 0.0513 & 0.1934 \\
cross\_train & F03-2 & vc & 811 & 951.9000 & 1073.4000 & 0.1184 & 0.0490 & 0.0393 & 0.2688 \\
cross\_train & F06-1 & vc & 122 & 1133.1000 & 1151.2500 & 0.1223 & 0.0437 & 0.0427 & 0.2175 \\
cross\_test & F03-4 & vc & 832 & 946.5000 & 1071.1500 & 0.0824 & 0.0341 & 0.0112 & 0.1699 \\
real\_f03\_gr & F03-4 & porosity & 1435 & 856.0500 & 1071.1500 & 0.2789 & 0.0142 & 0.2471 & 0.3264 \\
\bottomrule
\end{tabular}
}
\end{table}

The cross-well split is intentionally difficult. The held-out F03-4 Vc target
has a lower mean than the training wells, indicating a distribution shift that
direct supervised regressors must extrapolate across. This motivates a
generative constraint on plausible target windows rather than relying only on a
pointwise mapping from logs to targets. The low-data axis further increases this
difficulty by increasing the sliding-window step and reducing the number of
training windows.

\begin{figure}[htbp]
  \centering
  \includegraphics[width=0.86\linewidth]{eda_target_summary.png}
  \caption{Summary of target distributions used in the CLP-CSGM experiments.
  The held-out cross-well Vc target has a lower mean than the training wells,
  while the F03 porosity case has a narrower target range.}
  \label{fig:eda_target_summary}
\end{figure}

Figure~\ref{fig:eda_crosswell_distribution} shows the full Vc distributions for
the cross-well experiment. The held-out F03-4 distribution is shifted toward
lower Vc values relative to the training wells. This shift is important for the
article narrative: the method is not only interpolating within one well, but
using sparse target-well measurements to adapt a learned latent prior under a
visible distribution mismatch.

\begin{figure}[htbp]
  \centering
  \includegraphics[width=0.78\linewidth]{eda_crosswell_vc_distribution.png}
  \caption{Cross-well Vc target distributions for training wells and held-out
  F03-4. The held-out test distribution is shifted toward lower Vc, motivating a
  measurement-consistent method rather than a purely supervised predictor.}
  \label{fig:eda_crosswell_distribution}
\end{figure}

Table~\ref{tab:eda_correlations} reports simple Pearson correlations between
input channels and targets. Gamma ray is strongly correlated with Vc in both
cross-well train and test splits, whereas the real F03 porosity experiment uses
gamma ray only despite a weak linear correlation with porosity. This GR-only
choice is deliberate: it weakens the dense-input predictor and tests whether
sparse target-well measurements can improve a log-conditioned generative prior.

\begin{table}[htbp]
  \centering
  \caption{Channel-target Pearson correlations. These values are descriptive
  only and are not used as model inputs.}
  \label{tab:eda_correlations}
  \small
  \begin{tabular}{lllr}
\toprule
dataset & channel & target & pearson\_r \\
\midrule
cross\_train & sonic & vc & -0.1642 \\
cross\_train & rhob & vc & 0.0807 \\
cross\_train & gr & vc & 0.9733 \\
cross\_train & ai & vc & 0.1672 \\
cross\_train & vp & vc & 0.1819 \\
cross\_test & sonic & vc & -0.1606 \\
cross\_test & rhob & vc & -0.0061 \\
cross\_test & gr & vc & 0.9772 \\
cross\_test & ai & vc & 0.1330 \\
cross\_test & vp & vc & 0.1077 \\
real\_f03 & ac & porosity & 0.9996 \\
real\_f03 & gr & porosity & -0.1759 \\
\bottomrule
\end{tabular}

\end{table}

Figure~\ref{fig:eda_correlations} visualizes the same channel-target
correlations. In the cross-well Vc task, GR is the dominant linear correlate in
both train and held-out wells, which explains why direct log-based baselines are
not trivial competitors. In the F03 porosity task, however, GR is weakly and
negatively correlated with porosity, whereas AC is almost perfectly correlated.
Using GR only therefore creates a deliberately harder dense-input setting where
the sparse measurements and the latent prior have a clear role.

\begin{figure}[htbp]
  \centering
  \includegraphics[width=0.88\linewidth]{eda_channel_target_correlations.png}
  \caption{Channel-target Pearson correlations for the cross-well Vc and F03
  porosity data. GR is highly informative for Vc but weak for F03 porosity,
  motivating the GR-only porosity benchmark as a difficult target-well test.}
  \label{fig:eda_correlations}
\end{figure}

Figure~\ref{fig:eda_f03_profile} complements the correlation table by showing
the real F03-4 GR and porosity profiles along depth. The GR-only experiment is
therefore a deliberately weakened dense-input setting: porosity varies smoothly
over the interval, while GR alone is not a near-deterministic porosity proxy.
This makes the sparse porosity observations central to the reconstruction task.

\begin{figure}[htbp]
  \centering
  \includegraphics[width=0.82\linewidth]{eda_f03_gr_porosity_depth.png}
  \caption{Real F03-4 depth profiles used in the GR-only porosity benchmark.
  Gamma ray is the only dense input, while porosity is the target profile to be
  reconstructed from sparse observations and the conditional latent prior.}
  \label{fig:eda_f03_profile}
\end{figure}

The EDA therefore supports the experimental design. Cross-well Vc evaluates
low-data transfer under sparse target measurements, while the F03 GR-only case
tests whether CLP-CSGM can recover useful porosity structure when the dense
input alone is intentionally limited.
